\documentclass[11pt]{article}

\usepackage[margin=1in]{geometry}
\usepackage{amsmath,amssymb,amsthm}
\usepackage{graphicx}
\usepackage[hidelinks]{hyperref}

\newtheorem{theorem}{Theorem}
\newtheorem{lemma}{Lemma}
\newtheorem{proposition}{Proposition}
\newtheorem{remark}{Remark}

\newcommand{\htte}{\operatorname{ht}}
\newcommand{\ims}{\operatorname{ims}}
\newcommand{\cf}{\operatorname{cf}}
\newcommand{\rk}{\operatorname{rk}}
\newcommand{\Lev}{\operatorname{Lev}}
\newcommand{\SigmaProd}{\Sigma}

\begin{document}

\begin{center}
{\Large A positive answer to the height \(\omega_2\) Valdivia \(r\)-tree problem}

\medskip
{\small Run \texttt{fa\_banach\_001}; source arXiv:2305.11737}
\end{center}

\section*{Source problem}

Russo and Somaglia's paper \emph{Banach spaces of continuous functions without
norming Markushevich bases} proves that every scattered compact \(r\)-tree of
height strictly less than \(\omega_2\) is Valdivia. Immediately after that
result they ask the boundary question below.

\begin{figure}[h]
\centering
\includegraphics[width=\textwidth]{figures/open_problem_crop.png}
\caption{Russo--Somaglia, arXiv:2305.11737, page 15: Problem 5.9.}
\end{figure}

\noindent
Problem 5.9 asks:

\begin{quote}
Let \(T\) be a scattered compact \(r\)-tree with
\(\htte(T)=\omega_2\). Must \(T\) be Valdivia?
\end{quote}

\section*{Definitions and cited facts}

All trees are endowed with the coarse wedge topology. For \(t\in T\),
\[
V_t=\{u\in T:t\leq u\}, \qquad
\ims(t)=\{u\in T:t<u,\ \htte(u,T)=\htte(t,T)+1\}.
\]
A tree is an \(r\)-tree if \(|\ims(t)|<\omega\) whenever
\(\cf(t)\geq\omega_1\).  Put
\[
D(T)=\{t\in T:\cf(t)\leq\omega\}.
\]

We use three standard facts from the source and the cited compact-tree
literature.

\begin{itemize}
\item Russo--Somaglia Corollary 5.7: every scattered compact \(r\)-tree of
height \(<\omega_2\) is Valdivia \cite{RussoSomaglia}.
\item Russo--Somaglia Lemma 5.5: a scattered compact tree of height
\(\leq\omega_2\) can be decomposed at its top Cantor--Bendixson derivative
into pairwise disjoint clopen, \(\wedge\)-closed, lower-rank subtrees and a
finite union of chains of cardinality at most \(\omega_1\)
\cite{RussoSomaglia}.
\item Somaglia/Kalenda clopen-family criterion: a compact zero-dimensional
space, in particular a compact tree in the coarse wedge topology, is Valdivia
provided it has a dense set \(D\) and a \(T_0\)-separating family of clopen
sets that is point-countable on \(D\). For compact trees one may take
\(D=D(T)\) \cite{SomagliaTrees}.
\end{itemize}

\section*{Proof intuition}

The source paper's lower-height theorem fails to apply directly at
\(\omega_2\) because the known height induction reaches only levels below
\(\omega_2\).  The boundary case becomes accessible if one inducts instead on
Cantor--Bendixson rank.

Take a hypothetical counterexample of minimal Cantor--Bendixson rank.  The
source paper's scattered-tree decomposition cuts off all top-rank behavior
along finitely many chains.  Everything off those chains is a clopen subtree
of smaller Cantor--Bendixson rank, hence Valdivia by minimality, even if some
of those pieces still have height \(\omega_2\).  Collapse each such clopen
piece to a marker node.  The quotient is only a finite-spine tree, so its
height is strictly below \(\omega_2\), and it is Valdivia by Russo--Somaglia's
already-proved theorem.

The remaining point is a replacement lemma: a Valdivia zero-dimensional
quotient remains Valdivia after replacing marker nodes, lying in the dense
\(\Sigma\)-part, by clopen Valdivia compact fibers.  Pull back a
point-countable clopen separating family from the quotient and add the
point-countable separating families inside the fibers.

\section*{A clopen replacement lemma}

\begin{lemma}[Clopen replacement]\label{lem:replacement}
Let \(K\) and \(L\) be compact zero-dimensional spaces and let
\(q:K\to L\) be a continuous surjection.  Suppose that \(L\) is Valdivia,
witnessed by a dense set \(E\subset L\) and a \(T_0\)-separating family
\(\mathcal A\) of clopen subsets of \(L\) point-countable on \(E\).

Assume there is a set \(M\subset E\) such that:
\begin{enumerate}
\item if \(y\notin M\), then \(q^{-1}(y)\) is a singleton;
\item if \(m\in M\), then \(K_m=q^{-1}(m)\) is a clopen Valdivia compactum;
\item for every \(m\in M\), \(K_m\) has a dense set \(E_m\) and a
\(T_0\)-separating clopen family \(\mathcal A_m\), taken inside \(K_m\), that
is point-countable on \(E_m\).
\end{enumerate}
Finally assume that
\[
E_K=q^{-1}(E\setminus M)\cup\bigcup_{m\in M}E_m
\]
is dense in \(K\).
Then \(K\) is Valdivia.
\end{lemma}

\begin{proof}
For \(A\in\mathcal A\), the preimage \(q^{-1}(A)\) is clopen in \(K\).  Since
each \(K_m\) is clopen, every member of every internal family
\(\mathcal A_m\) is also clopen in \(K\).  Define
\[
\mathcal B=\{q^{-1}(A):A\in\mathcal A\}
       \cup \bigcup_{m\in M}\mathcal A_m
\]
and
\[
E_K=q^{-1}(E\setminus M)\cup\bigcup_{m\in M}E_m.
\]
The family \(\mathcal B\) is \(T_0\)-separating.  If two points have different
images in \(L\), a member of \(\mathcal A\) separates their images, and its
preimage separates the points.  If they have the same image, then they lie in
one fiber \(K_m\), and a member of \(\mathcal A_m\) separates them.

Finally, \(\mathcal B\) is point-countable on \(E_K\).  If
\(x\in q^{-1}(E\setminus M)\), the only members of \(\mathcal B\) containing
\(x\) are pullbacks of members of \(\mathcal A\) containing \(q(x)\), a
countable family.  If \(x\in E_m\), then the pullback part contributes only
the countably many \(A\in\mathcal A\) with \(m\in A\), and the internal part
contributes only the countably many members of \(\mathcal A_m\) containing
\(x\).  Therefore \(K\) has a dense set with a point-countable
\(T_0\)-separating clopen family, so \(K\) is Valdivia by the criterion.
\end{proof}

\section*{The finite-spine quotient}

\begin{lemma}[Finite-spine quotient]\label{lem:quotient}
Let \(T\) be a scattered compact \(r\)-tree with \(\htte(T)\leq\omega_2\) and
Cantor--Bendixson rank \(\rho+1\).  There is a continuous surjection
\[
q:T\longrightarrow Q
\]
onto a scattered compact \(r\)-tree \(Q\) of height \(<\omega_2\) such that:
\begin{enumerate}
\item every non-singleton fiber of \(q\) is a clopen, \(\wedge\)-closed,
compact subtree of \(T\) of Cantor--Bendixson rank \(<\rho+1\);
\item the points of \(Q\) corresponding to non-singleton fibers are on
successor levels, or are the root of \(Q\), hence belong to \(D(Q)\);
\item singleton fibers are exactly the points lying on a finite union of
chains of \(T\).
\end{enumerate}
\end{lemma}

\begin{proof}
Russo--Somaglia Lemma 5.5 gives the required decomposition.  We recall the
part of that construction that matters here.  Since \(T\) is compact and
scattered, \(T^{(\rho)}\) is finite.  Around each \(t\in T^{(\rho)}\) choose
a clopen wedge neighborhood \(W_{s_0(t)}^{F_t}\) isolating \(t\) from the
other top-rank points.  The proof of Lemma 5.5 then writes \(T\) as a
pairwise disjoint union of:
\begin{itemize}
\item a clopen \(\wedge\)-closed remainder
\[
S=T\setminus \bigcup_{t\in T^{(\rho)}}V_{s_0(t)}
\]
with rank \(<\rho+1\);
\item clopen \(\wedge\)-closed side pieces \(V_r\) occurring along the finite
chains from \(s_0(t)\) to \(t\), each of rank \(<\rho+1\);
\item clopen \(\wedge\)-closed pieces \(W_r^{G_r}\) immediately above the
top-rank points, again each of rank \(<\rho+1\);
\item the finite union of residual chains
\[
C=\bigcup_{t\in T^{(\rho)}}\bigl(\hat t\cap V_{s_0(t)}\bigr).
\]
\end{itemize}
The rank inequalities are exactly the point of the decomposition in
Russo--Somaglia's proof of Theorem 5.6.

Collapse every clopen lower-rank piece in the first three classes to one
marker point and leave the residual chain points distinct.  The order on the
quotient \(Q\) is the evident tree order: a marker replacing a side piece
takes the position of the root of that side piece; the marker replacing
\(S\), if \(S\neq\emptyset\), is the root below the finitely many residual
spines.  Because the residual part is a finite union of chains and every
point of \(T\) has height \(<\omega_2\), the quotient has height
\(<\omega_2\).  Adding marker successors cannot reach height \(\omega_2\).

The quotient is chain complete and rooted: every chain either is eventually
inside one residual chain, or terminates at a marker replacing one collapsed
clopen subtree.  The \(r\)-tree condition is inherited from \(T\).  At a
quotient point of cofinality at least \(\omega_1\), its immediate successors
come from immediate successors of the corresponding point of \(T\), and there
are finitely many because \(T\) is an \(r\)-tree.  The quotient is scattered
because it is a finite union of ordinal chains with isolated marker
successors.

Finally, \(q\) is continuous.  The basic clopen subsets of \(Q\) are generated
by wedges \(V_a\), where \(a\) is the root or lies on a successor level, and
their complements.  If \(a\) is a marker, \(q^{-1}(V_a)\) is precisely the
corresponding clopen piece of \(T\).  If \(a\) lies on a residual chain, then
\(q^{-1}(V_a)\) is the corresponding wedge in \(T\), possibly with finitely
many already clopen collapsed side pieces removed or added according to the
finite Boolean description of the basic wedge.  Hence each basic clopen of
\(Q\) has clopen preimage in \(T\).  This proves the asserted properties.
\end{proof}

\section*{Main theorem}

\begin{theorem}\label{thm:main}
Let \(T\) be a scattered compact \(r\)-tree with \(\htte(T)=\omega_2\).  Then
\(T\) is Valdivia.
\end{theorem}

\begin{proof}
Assume, toward a contradiction, that there is a non-Valdivia scattered
compact \(r\)-tree of height \(\omega_2\).  Choose such a tree \(T\) with
minimal Cantor--Bendixson rank.

Apply Lemma \ref{lem:quotient}.  For every non-singleton fiber
\(P=q^{-1}(m)\), the set \(P\) is a clopen, \(\wedge\)-closed, compact
subtree of \(T\), and its Cantor--Bendixson rank is strictly smaller than the
rank of \(T\).  It is again scattered and an \(r\)-tree, with height at most
\(\omega_2\).  If \(\htte(P)<\omega_2\), then \(P\) is Valdivia by
Russo--Somaglia Corollary 5.7.  If \(\htte(P)=\omega_2\), then \(P\) is
Valdivia by the minimality of \(T\)'s Cantor--Bendixson rank.  Thus every
non-singleton fiber of \(q\) is Valdivia.

The quotient \(Q\) is a scattered compact \(r\)-tree of height \(<\omega_2\),
so \(Q\) is Valdivia by Russo--Somaglia Corollary 5.7.  Moreover, the marker
points corresponding to non-singleton fibers are in \(D(Q)\), by
Lemma \ref{lem:quotient}.  For each non-singleton fiber \(P\), choose a dense
\(\Sigma\)-set and a point-countable clopen separating family witnessing that
\(P\) is Valdivia.  For \(Q\), choose the corresponding witness with
\(D(Q)\) as the dense set, as supplied by the compact-tree criterion.
The dense set required in Lemma \ref{lem:replacement} is dense in \(T\):
inside each clopen fiber this is the density of its chosen dense
\(\Sigma\)-set, and along the residual finite spine it is the usual density
of points of countable cofinality in a compact tree.

All hypotheses of the clopen replacement Lemma \ref{lem:replacement} are now
satisfied.  Therefore \(T\) is Valdivia, contradicting the choice of \(T\).
This proves the theorem.
\end{proof}

\section*{Relation to the earlier partial packet}

The earlier root-split packet handled the special case where every immediate
successor subtree above the root has height \(<\omega_2\).  The proof above
removes that extra hypothesis.  Subtrees of height \(\omega_2\) may occur
below the root, but once the tree is split by Cantor--Bendixson rank, those
tall subtrees have lower rank and are handled by the minimal-rank induction.

\section*{Verification and novelty notes}

The proof is non-computational.  The local verification script checks that the
packet includes the source problem, the full theorem, the replacement lemma,
and the minimal-rank argument.

Local cheap-index searches covered the arXiv id \texttt{2305.11737}, source
title, \texttt{Problem 5.9}, \texttt{scattered compact r-tree},
\texttt{height omega\_2}, \texttt{Valdivia compact tree}, and nearby
norming-M-basis phrases.  The only existing packet for this arXiv id was the
root-split partial subcase produced earlier in this run.

External web searches on June 16, 2026 used the exact source title, arXiv id,
\texttt{Can the first statement of Theorem 3.2 be reversed}, and variants of
\texttt{compact r-tree height omega\_2 Valdivia}.  They found the source paper
and Somaglia's cited compact-tree characterization, but no later paper
explicitly resolving Problem 5.9.

\section*{Human review recommendation}

Human review should focus first on Lemma \ref{lem:quotient}.  In particular,
check that the quotient obtained from the Russo--Somaglia Lemma 5.5
decomposition is indeed a compact \(r\)-tree of height \(<\omega_2\), and that
the quotient map pulls back basic clopen sets to clopen sets of \(T\).  Once
that finite-spine quotient is accepted, Lemma \ref{lem:replacement} and the
minimal Cantor--Bendixson rank contradiction are routine.

\begin{thebibliography}{9}

\bibitem{RussoSomaglia}
T. Russo and J. Somaglia,
\emph{Banach spaces of continuous functions without norming Markushevich
bases},
arXiv:2305.11737.

\bibitem{SomagliaTrees}
J. Somaglia,
\emph{On compact trees with the coarse wedge topology},
arXiv:1803.11107.

\bibitem{Kalenda}
O. F. K. Kalenda,
\emph{Valdivia compact spaces in topology and Banach space theory},
Extracta Math. 15 (2000), 1--85.

\end{thebibliography}

\end{document}
